\section{Introducción}
Hoy en día, una gran parte del código que usan las empresas es ya viejo. A este lío se le conoce como código legacy, y es un dolor de cabeza que no para de crecer. Nace cuando, por la prisa, se prioriza terminar una función rápido en vez de hacerla con calidad. Esto nos deja sistemas rígidos que son carísimos de mantener y que, a la larga, aumentan los tiempos de desarrollo y la cantidad de errores (bugs).

El problema es grande: un sistema así de tieso no solo complica la vida del programador, sino que es lento para innovar y afecta directamente la seguridad. Si el software es inflexible, no puede crecer. Por eso, este artículo no se queda en la teoría, sino que presenta una salida real y sostenible para la evolución del software, tomando como referencia el desarrollo del Portal Agro-comercial del Huila.

Entonces, en lugar de creer que hay que tirar todo a la basura y empezar de cero, lo que demostramos es que sí se puede refactorizar. Y no solo el código, sino la forma en que trabajamos, transformando esa famosa deuda técnica en una verdadera ventaja. Para conseguir esto, nos enfocamos en ideas muy concretas: seguir reglas de diseño fuertes (como SOLID), asegurar la calidad con cosas sencillas como las pruebas automáticas (TDD) y trabajar de forma ágil con Scrum. Todo esto se aterrizó en una arquitectura modular, usando herramientas como C\# ASP.NET Core y Angular, justo como lo documentamos en nuestro proyecto.

Para que puedas entender todo el proceso, el artículo sigue este recorrido:

En la Sección 2, primero te explicaremos la teoría que sustenta todo (el marco conceptual).

Después, en la Sección 3, detallaremos las metodologías que realmente aplicamos en el día a día.

En la Sección 4 contaremos cómo fue la implementación real en el Portal Agro-comercial del Huila.

Las Secciones 5, 6 y 7 están dedicadas a los resultados, la discusión de las lecciones aprendidas y las conclusiones finales.

Nuestra contribución más importante es doble: un marco teórico unificado y la demostración práctica de su éxito con un caso de estudio real, lo que se convierte en una guía práctica para otros equipos de ingeniería. El impacto esperado es claro: tener sistemas mucho más fáciles de mantener y generar un cambio en el equipo hacia una mentalidad de calidad continua.